% Options for packages loaded elsewhere
\PassOptionsToPackage{unicode}{hyperref}
\PassOptionsToPackage{hyphens}{url}
\PassOptionsToPackage{dvipsnames,svgnames,x11names}{xcolor}
%
\documentclass[
  11pt,
]{article}
\usepackage{amsmath,amssymb}
\usepackage{iftex}
\ifPDFTeX
  \usepackage[T1]{fontenc}
  \usepackage[utf8]{inputenc}
  \usepackage{textcomp} % provide euro and other symbols
\else % if luatex or xetex
  \usepackage{unicode-math} % this also loads fontspec
  \defaultfontfeatures{Scale=MatchLowercase}
  \defaultfontfeatures[\rmfamily]{Ligatures=TeX,Scale=1}
\fi
\usepackage{lmodern}
\ifPDFTeX\else
  % xetex/luatex font selection
\fi
% Use upquote if available, for straight quotes in verbatim environments
\IfFileExists{upquote.sty}{\usepackage{upquote}}{}
\IfFileExists{microtype.sty}{% use microtype if available
  \usepackage[]{microtype}
  \UseMicrotypeSet[protrusion]{basicmath} % disable protrusion for tt fonts
}{}
\makeatletter
\@ifundefined{KOMAClassName}{% if non-KOMA class
  \IfFileExists{parskip.sty}{%
    \usepackage{parskip}
  }{% else
    \setlength{\parindent}{0pt}
    \setlength{\parskip}{6pt plus 2pt minus 1pt}}
}{% if KOMA class
  \KOMAoptions{parskip=half}}
\makeatother
\usepackage{xcolor}
\usepackage[margin=1in]{geometry}
\setlength{\emergencystretch}{3em} % prevent overfull lines
\providecommand{\tightlist}{%
  \setlength{\itemsep}{0pt}\setlength{\parskip}{0pt}}
\setcounter{secnumdepth}{-\maxdimen} % remove section numbering
\ifLuaTeX
  \usepackage{selnolig}  % disable illegal ligatures
\fi
\usepackage{bookmark}
\IfFileExists{xurl.sty}{\usepackage{xurl}}{} % add URL line breaks if available
\urlstyle{same}
\hypersetup{
  pdftitle={A nodal compatibility obstruction for finite-energy equations},
  pdfauthor={Local research continuation for P16.1},
  colorlinks=true,
  linkcolor={Maroon},
  filecolor={Maroon},
  citecolor={Blue},
  urlcolor={Blue},
  pdfcreator={LaTeX via pandoc}}

\title{A nodal compatibility obstruction for finite-energy equations}
\author{Local research continuation for P16.1}
\date{4 October 2026}

\begin{document}
\maketitle

Fresh root-authored local diagnostic. It has not been sent to either
running native author, independently reviewed, or used to close P16.1.
The bounded-metric obstruction below is complete within its stated
endpoint. The finite-energy consequence states its exact interfaces
rather than assuming an unspecified singular-space theory.

\subsection{1. The algebraic curve and its actual
polarization}\label{the-algebraic-curve-and-its-actual-polarization}

Take \(K=\mathbb C((t))\) with \(0<|t|<1\). Let \(D\) be an integral
rational projective nodal curve, with normalization
\(\nu:\mathbb P^1_K\to D\) identifying the rational points \(0\) and
\(\infty\) to its sole singular point \(q\). Such a curve is obtained
from a rational plane nodal cubic by a projective change of the
normalization parameter. It is geometrically integral and proper; it is
not geometrically unibranch at \(q\).

Glue the fibers of \(\mathcal O_{\mathbb P^1}(2)\) at \(0\) and
\(\infty\) using the identity in the standard local frames. This gives a
line bundle \(L\) on \(D\) with \(\nu^*L=\mathcal O(2)\). Locally at the
node its sections are pairs of regular functions with the same value at
the two branches. This module is free of rank one over the nodal local
ring after choosing the matching frames, so the gluing gives an
invertible sheaf, rather than only a torsion-free rank-one sheaf. Its
degree is two; on an integral projective curve positive degree is ample.

In the affine normalization coordinate \(T\), the sections

\[
s_0=1+T^2,\qquad s_1=T
\]

of \(\mathcal O(2)\) satisfy the gluing condition: the constant
coefficient equals the coefficient of \(T^2\). They descend to \(L\) and
have no common zero, including at the node. Their Fubini--Study metric
\(\psi\) is a genuine global semipositive metric on \((D,L)\). Its
pullback has local weight

\[
\psi^\nu(T)=\log\max(|1+T^2|,|T|)=2\log\max(1,|T|).
\]

The equality holds at every Berkovich point. For \(|T|<1\) the first
term has norm one; for \(|T|>1\) it has norm \(|T|^2\); for \(|T|=1\)
the second term has norm one, so cancellation in \(1+T^2\) cannot change
the maximum. Hence the raw curvature measure upstairs is
\(2\delta_{\zeta_1}\), where \(\zeta_s\) is the radius-\(s\) Gauss
point. The raw degree and reference mass are two, not one.

\subsection{2. A prescribed divisorial measure with a non-descending
solution}\label{a-prescribed-divisorial-measure-with-a-non-descending-solution}

Choose \(r=|t|<1\) and set

\[
\mu_r=2\delta_{\nu(\zeta_r)}.
\]

This positive Radon measure has raw mass \(\deg L=2\), is supported at a
divisorial point in the nonsingular locus, and charges neither the node
nor any polar rigid point. On the smooth normalization the semipositive
metric

\[
\phi_r(T)=2\log\max(r,|T|)
\]

on \(\mathcal O(2)\) has curvature \(2\delta_{\zeta_r}\). Its difference
from the reference is the finite continuous function

\[
u_r=2\log\frac{\max(r,|T|)}{\max(1,|T|)}.
\]

It equals \(2\log r\) on \(|T|\leq r\), equals \(2\log|T|\) on
\(r\leq|T|\leq1\), and equals zero on \(|T|\geq1\), including at
infinity. In particular,

\[
u_r(0)=2\log r\ne0=u_r(\infty).
\]

For an actual metric on \(L\), its difference from the fixed descended
reference is one function on \(D^{\mathrm{an}}\). Its pullback must have
the SAME value at \(0\) and \(\infty\). Adding a global constant does
not remove the displayed discrepancy.

\textbf{Bounded endpoint obstruction.} There is no continuous global
semipositive metric on \(L\) whose genuine Chambert--Loir measure is
\(\mu_r\).

\textbf{Proof.} For a continuous semipositive metric \(\phi\) on \(L\),
normalization pullback is a continuous semipositive metric on
\(\mathcal O(2)\). The finite birational projection formula gives

\[
\nu_*\operatorname{MA}_{\mathbb P^1}(\nu^*\phi)
=\operatorname{MA}_D(\phi).
\]

For model metrics this is the degree-one intersection projection formula
on a common model. For continuous semipositive metrics take their
uniform model/FS approximations; pullback preserves the same uniform
error and both measures converge weakly. Equivalently, the
smooth-curvature version follows from
\(\nu_*\delta_{\mathbb P^1}=\delta_D\) and the projection formula for
the analytic fundamental class, then extends uniformly. The standard
continuous endpoint over general projective geometrically reduced
schemes is available in Boucksom--Eriksson, Theorem 8.4, and its
birational energy projection formula is written in Proposition
9.14(vii). This application is a curvature-measure projection formula,
not merely an inference from an energy equality without the projection
argument.

If the downstairs measure were \(\mu_r\), positivity and that formula
force the entire upstairs measure to be \(2\delta_{\zeta_r}\): the
measure of the node is zero, hence no upstairs mass lies at its two
preimages; everywhere else normalization is an isomorphism. Smooth-curve
uniqueness for positive Radon measures, as recalled and proved by the
current/Green-function argument in BFJ, Section 9.1, gives
\(\nu^*\phi=\phi_r+c\). The difference from \(\psi^\nu\) then has
unequal values at \(0,\infty\), contradicting metric descent. This
proves the obstruction.

\subsection{3. The measure is not excluded by the usual dual-energy
test}\label{the-measure-is-not-excluded-by-the-usual-dual-energy-test}

Write \(\ell=-\log r>0\). On the smooth normalization use the
probability-normalized energy relative to \(\psi^\nu\),

\[
\mathcal E_{\mathbb P^1}(u)
=\frac1{2\deg L}\int u\,
 \bigl(\operatorname{MA}(\psi^\nu+u)+\operatorname{MA}(\psi^\nu)\bigr).
\]

This normalization is only for the numerical dual-energy condition; the
equation and its measure still have raw mass two. With \(\deg L=2\),
integration by parts gives

\[
\mathcal E_{\mathbb P^1}(u)=u(\zeta_1)-\tfrac14\mathcal D(u),
\qquad \mathcal D(u)=-\int u\,dd^c u.
\]

For finite-energy potentials the same formula is the Dirichlet-Hilbert
extension. Thuillier Section 3.5 proves the Hilbert realization, bounded
evaluation at nonrigid points and the dual-energy criterion; Proposition
3.5.9 is stated with the exact sign
\(2\int\varphi\,d\nu-\mathcal D(\varphi)\) for a zero-mass measure
\(\nu\). Its actual PDF pages 118--119 were inspected to resolve OCR
ambiguity. We use the following elementary special case of that
framework, with constants computed explicitly.

The segment \([\zeta_r,\zeta_1]\) has metric length \(\ell\). Restrict a
finite-graph function \(u\) to that segment. The fundamental theorem of
calculus and Cauchy--Schwarz give

\[
|u(\zeta_1)-u(\zeta_r)|^2\leq\ell\,\mathcal D(u).
\]

Here the segment contribution is at most the full nonnegative Dirichlet
energy. Completion in Thuillier's \(W^1\), together with continuity of
the two nonrigid evaluations, proves the same bound for every
normalization finite-energy potential. Consequently, with
\(z=\sqrt{\mathcal D(u)}\),

\[
\mathcal E_{\mathbb P^1}(u)-u(\zeta_r)
\leq\sqrt\ell\,z-\tfrac14z^2\leq\ell.
\]

The last bound is the completed square \(\ell-(z/2-\sqrt\ell)^2\). For
\(u=u_r\), \(\mathcal D(u_r)=4\ell\), \(\mathcal E(u_r)=-\ell\) and
\(u_r(\zeta_r)=-2\ell\), so equality holds upstairs. No infinity-valued
point evaluation is used in this energy estimate.

For the continuous metric class on \(D\), normalized birational energy
invariance gives

\[
\mathcal E_D(\phi-\psi)=
\mathcal E_{\mathbb P^1}\bigl(\nu^*(\phi-\psi)\bigr).
\]

Also \(\int(\phi-\psi)\,d\mu_r/2\) is its pullback's value at
\(\zeta_r\). Thus the dual energy of \(\mu_r/2\) on that class is
between zero and \(\ell\): zero is a lower bound from the reference
metric, and the displayed inequality gives the upper bound.

The same conclusion holds on any monotone finite-energy completion for
which pullback and the energy are the decreasing limits of their
continuous counterparts. Indeed, the continuous birational equality then
passes to the defining limits and the pullback has finite normalization
energy. Its value at the nonrigid \(\zeta_r\) is finite by the Hilbert
evaluation result, and the same estimate applies. Thus this concrete
target has finite dual energy for that exact completion. This does not
identify an unspecified larger singular-space energy cone with that
completion.

\subsection{4. Precisely stated finite-energy
consequence}\label{precisely-stated-finite-energy-consequence}

Consider a proposed equation theory on this \((D,L)\) with the following
interfaces:

\begin{enumerate}
\def\labelenumi{\arabic{enumi}.}
\tightlist
\item
  Its admissible potentials are actual psh metrics on \(L\), so
  normalization pulls them back to the usual psh metric class on
  \(\mathcal O(2)\), with the values of their reference differences
  equal at the two normalization preimages of \(q\).
\item
  Its finite-energy domain is the usual monotone extension of the
  continuous model energy, with the pullback/energy compatibility
  described in Section 3. In particular the target \(\mu_r\) meets its
  raw mass and finite dual-energy conditions.
\item
  Its measure operator agrees on the nonsingular locus with the usual
  one-dimensional curvature/Monge--Ampere operator and is positive of
  raw mass two.
\end{enumerate}

\textbf{Consequence.} Such a theory cannot have an existence theorem for
every measure of mass two and finite dual energy.

\textbf{Proof.} A solution for \(\mu_r\) would pull back to a psh metric
\(\phi^\nu\) on \(\mathcal O(2)\) of finite energy. On
\(\mathbb P^1\setminus\{0,\infty\}\), locality and interface 3 force its
curvature to be \(2\delta_{\zeta_r}\). Its usual positive global
curvature has total mass two. This is the degree formula on the smooth
normalization; for finite energy its full-mass measure also has total
mass two. There is therefore no remaining positive mass at \(0\) or
\(\infty\). The normalization equation is exactly the one solved by
\(\phi_r\). Curve uniqueness gives \(\phi^\nu=\phi_r+c\), again
contradicting interface 1. This proof only needs locality on the
normalization-isomorphic locus and the known smooth-curve mass; it does
not assume an unproved global finite-energy projection formula on the
singular curve.

The example explains why merely normalizing a multibranch curve cannot
prove the requested equation theorem. It satisfies projectivity,
ampleness, discrete valuation, characteristic-zero residue and
divisorial support; those axes cause no failure here. The conflict is
the conductor's equality condition on metric values. A geometrically
unibranch curve has no pair of preimages with this condition, so this
example is not an obstruction in that range. Higher-dimensional singular
and the genuinely unibranch extension of original P16.1 remain separate
obligations.

\subsection{5. Sources, status and exact
limits}\label{sources-status-and-exact-limits}

The published primary sources actually read for this diagnostic are: BFJ
arXiv:1201.0188v1 Section 9.1 (literal extract LF1829--1875);
Boucksom--Eriksson arXiv:1805.01016v3 (LF2620--2708 for the continuous
endpoint, LF3540--3588 for energy and birational invariance);
Pille--Schneider arXiv:2209.04879v2 (LF1932--1997, including the
distinction between the continuous operator and general finite energy);
and Thuillier tel-00010990v1 (LF5793--5905 and LF6067--6115, with actual
PDF pages 114,118,119 inspected). Standard algebraic facts about a
rational nodal curve, gluing invertible sheaves and positive-degree
ampleness are explicitly used. No unpublished generated result is a
premise.

The continuous metric obstruction and the conditional finite-energy
consequence have exact but different scopes. The latter does not prove
that every conceivable singular-space positivity convention has the
three interfaces. It does show that an existence theorem with these
usual metric, energy and locality signatures is incompatible with this
prescribed measure. There is no independent mathematical acceptance yet
and no original P16.1 closure. The active authors' inputs and work have
not been amended or influenced.

\end{document}
